% Legislative Findings and Evidentiary Record. Numbered-finding legislative
% style, three groups, one External-Standards table (Group B), and one overall
% evidence table. The dense numbers live in the tables.

The Congress finds the following. The findings below are grouped to track the
three pillars of the record on which this bill rests: that large language models
are appropriate for verification of robot-patient code, that the verification can
be generated and executed today against published external standards, and that
the readiness standards and simulation record make the robot-patient case
concrete. Each finding states a legislative conclusion and points, by reference
to Table~\ref{tab:findings-evidence}, to the recorded evidence that supports it;
the external standards that anchor the second pillar are enumerated in
Table~\ref{tab:findings-standards}.

\subsection{Group A: Large Language Models Are Appropriate for Code Verification}

\textbf{Finding A1.} An autonomous large language model agent can build, execute,
and document a complete verification, validation, and uncertainty quantification
pipeline under strict assurance, and when it does, the verification gate, not the
generation step, is the decision-bearing work. In the author's first study the
pipeline emitted a complete deliverable in under a millisecond, while the gate
accepted only one of six candidate artifacts, blocked five, and escalated one,
holding the verification fraction at exactly 1.0; the supporting suite passed 51
of 51 tests \cite{kawchak2026vvuq01}.

\textbf{Finding A2.} The same priority holds when the system is an embodied
surgical humanoid rather than an analytics pipeline. In the author's second study
the assurance suite carried 64 of 172 tests, more than a third of the entire test
budget, which is itself evidence that verifying the behavior outweighs generating
it \cite{kawchak2026vvuq02}. Findings A1 and A2 together establish that the gate
is the substantial work and therefore the proper object of a statutory mandate.

\subsection{Group B: Practical VVUQ Against External Standards}

\textbf{Finding B1.} Verification, validation, and uncertainty quantification of
robot-patient interaction code are practical today: the author's second study
generated the control software for a surgical humanoid and executed a ten-gate
assurance suite over it, without manual intervention and reproducibly from a
fixed seed \cite{kawchak2026vvuq02}.

\textbf{Finding B2.} Crucially, every gate was bound to published consensus
standards already used in real life, not to ad hoc thresholds chosen for
convenience. This is the feature that makes the assurance argument traceable and
defensible to a regulator: a reviewer can follow each numeric limit back to the
standard that sets it. The fourteen external standards and two clinical baselines,
and the role each plays at the gate, are enumerated in
Table~\ref{tab:findings-standards}; their per-gate bindings are recorded in the
standards map referenced in Table~\ref{tab:findings-evidence}
\cite{asmevv40, nasastd7009a, iec8060127, iec60601, isots15066, iso13482,
iso10218, iso9283, iec62304, iso14971, iso13849, ul4600, ieee7009,
fda2023credibility}.

\textbf{Finding B3.} The second study produced two artifacts that demonstrate the
method at scale: a 1{,}790-line four-entrant comparison and a 1{,}001-row
non-repetitive humanoid sensor stream, each regenerated byte for byte from seed
20260525. These artifacts are referenced in Table~\ref{tab:findings-evidence} and
treated in detail in Sections~\ref{sec:algorithm_documentation}
and~\ref{sec:supporting_documentation}; they are named here only as evidence that
the generation-and-execution loop yields auditable records rather than claims.

\begin{table}[ht]
\centering
\begin{tabularx}{\textwidth}{L{3.5cm} Y L{1.4cm}}
\toprule
\textbf{External standard} & \textbf{Role at the gate} & \textbf{Source} \\
\midrule
\multicolumn{3}{l}{\textit{Cross-cutting credibility and software-process standards}} \\
ASME V\&V 40-2018 & Model credibility framework: context of use, model risk, credibility goals & \cite{asmevv40} \\
NASA-STD-7009A & Models-and-simulations credibility; uncertainty quantification as a first-class factor & \cite{nasastd7009a} \\
FDA CM\&S credibility (2023) & Regulatory acceptance of the V\&V 40 credibility argument & \cite{fda2023credibility} \\
IEC 62304 & Medical-device software life cycle; verification posture and safety class & \cite{iec62304} \\
ISO 14971:2019 & Risk management: model risk, risk-control verification, recorded human review & \cite{iso14971} \\
\multicolumn{3}{l}{\textit{Robotic-surgery and medical-electrical standards}} \\
IEC 80601-2-77:2019 & Basic safety and essential performance of robotically assisted surgery & \cite{iec8060127} \\
IEC 60601-1 & Single-fault safe state for the balance and fault e-stop paths & \cite{iec60601} \\
ISO 9283:1998 & Pose accuracy and repeatability for fingertip placement and handover & \cite{iso9283} \\
\multicolumn{3}{l}{\textit{Collaborative-, service-, and industrial-robot safety standards}} \\
ISO/TS 15066:2016 & Biomechanical contact-force and clearance limits & \cite{isots15066} \\
ISO 13482:2014 & Stability and fall prevention for a standing humanoid & \cite{iso13482} \\
ISO 10218-1:2011 & Speed and separation monitoring for shared-operating-room collision & \cite{iso10218} \\
ISO 13849-1:2023 & Performance level of the fault and e-stop control path & \cite{iso13849} \\
\multicolumn{3}{l}{\textit{Autonomy safety standards}} \\
UL 4600 & Autonomy safety case and safe-state default & \cite{ul4600} \\
IEEE 7009-2024 & Fail-safe default to a safe state under fault or uncertainty & \cite{ieee7009} \\
\multicolumn{3}{l}{\textit{Clinical baselines}} \\
Dutch human-surgeon cohort & Human-outcome reference for the four-entrant comparison & \cite{kawchak2026vvuq02} \\
Callery Fistula Risk Score & Pancreatic-fistula reference binding the anastomosis gate & \cite{kawchak2026vvuq02} \\
\bottomrule
\end{tabularx}
\caption{The fourteen external standards and two clinical baselines that bind the VVUQ gates.}
\label{tab:findings-standards}
\end{table}

\subsection{Group C: Two Author Standards and Five Simulations}

\textbf{Finding C1.} Physical AI VVUQ is supported by two author readiness
standards, the Unification Standard Level for robots and the Physical AI Standard
Level for sites, which together make the robot-patient case concrete by fixing
what readiness means and how it is measured \cite{kawchak2026national,
kawchak2026sitedocs}.

\textbf{Finding C2.} The robot-patient case is further supported by five trial,
sponsor, and mobile simulations, summarised with one headline metric each in
Table~\ref{tab:findings-evidence}: a 24-hour multi-patient trial simulation, a
single-patient ten-stage journey, a fully automated sponsor simulation, a mobile
pancreatic-cancer surgical-humanoid simulation, and an accelerated
patient-prediction study of four extensive simulations \cite{kawchak2026sponsor,
kawchak2026vvuq02, kawchak2026prediction}. A federated-learning framework with a
triple artificial-intelligence peer review of model updates before deployment
provides the supporting multi-site infrastructure \cite{kawchak2026federated}.

\begin{table}[ht]
\centering
\begin{tabularx}{\textwidth}{L{2.0cm} L{5.2cm} Y}
\toprule
\textbf{Finding} & \textbf{Source record (abbreviated)} & \textbf{Headline metric} \\
\midrule
\multicolumn{3}{l}{\textit{Parent: papers/VVUQ-01/ and papers/VVUQ-02/}} \\
A1 & VVUQ-01 execution/README.md & 51/51 tests; 1 ACCEPT, 5 BLOCK, 1 ESCALATE \\
A2 & VVUQ-02 execution/README.md & 172 tests; 64 in the 10-gate suite \\
B1--B3 & VVUQ-02 codegen/docs/standards\_map.md & 14 standards + 2 clinical baselines, per gate \\
\multicolumn{3}{l}{\textit{Parent: physical-ai .../final\_paper/sections/ and Zenodo}} \\
C1 & psl\_usl\_standards.tex; site docs & USL 1.0--10.0; PSL pass/fail (site) \\
C2a & site\_establishment.tex & 168 patients, 29 robots, 99.7\% uptime, 0 harm \\
C2b & patient\_journey.tex & 1 patient, 10 stages, 1{,}120 days (Stage IIIB NSCLC) \\
C2c--C2e & sponsor; H2 humanoid; prediction (Zenodo) & 288 decisions; mobile Whipple; 4 simulations \\
\bottomrule
\end{tabularx}
\caption{Evidentiary record underlying the legislative findings.}
\label{tab:findings-evidence}
\end{table}
